\documentclass[UTF8,fontset=fandol,openany,zihao=-4]{ctexbook}
% Shared lecture-note preamble for Big Data and Management Decision-Making.
% Chapter files follow latex_config.md and remain independently compilable.

\usepackage[a4paper,margin=2.6cm]{geometry}
\usepackage{amsmath,amssymb,amsthm,mathtools,bm}
\usepackage{xcolor}
\usepackage[most]{tcolorbox}
\usepackage{booktabs,longtable,array}
\usepackage{enumitem}
\usepackage{hyperref}
\usepackage{listings}
\usepackage{caption}
\usepackage{tikz}
\usetikzlibrary{arrows.meta,positioning,shapes.geometric}

\newcommand{\BDMFandolPath}{/usr/local/texlive/2025/texmf-dist/fonts/opentype/public/fandol/}
\setCJKmainfont[
  Path=\BDMFandolPath,
  UprightFont=FandolSong-Regular.otf,
  BoldFont=FandolSong-Regular.otf,
  ItalicFont=FandolKai-Regular.otf,
  BoldItalicFont=FandolKai-Regular.otf
]{FandolSong-Regular.otf}
\setCJKsansfont[
  Path=\BDMFandolPath,
  UprightFont=FandolSong-Regular.otf,
  BoldFont=FandolSong-Regular.otf
]{FandolSong-Regular.otf}
\setCJKmonofont[
  Path=\BDMFandolPath,
  UprightFont=FandolSong-Regular.otf,
  BoldFont=FandolSong-Regular.otf
]{FandolSong-Regular.otf}
\setmonofont{Menlo}
\linespread{1.13}

\hypersetup{
  colorlinks=true,
  linkcolor=blue!60!black,
  citecolor=blue!60!black,
  urlcolor=blue!60!black,
  pdfauthor={大数据与管理决策基础},
  pdfsubject={大数据与管理决策基础课程讲义}
}

\ctexset{
  chapter = {
    format = \huge\mdseries,
    name = {第,讲},
    number = \arabic{chapter},
    aftername = \quad
  },
  section = {format = \Large\mdseries},
  subsection = {format = \large\mdseries},
  subsubsection = {format = \normalsize\mdseries}
}

\setlist[itemize]{leftmargin=2em,itemsep=0.25em,topsep=0.25em}
\setlist[enumerate]{leftmargin=2em,itemsep=0.25em,topsep=0.25em}

\definecolor{BDMBlue}{RGB}{22,44,73}
\definecolor{BDMTeal}{RGB}{22,119,119}
\definecolor{BDMGray}{RGB}{76,84,95}
\definecolor{BDMLight}{RGB}{245,247,250}
\definecolor{BDMAmber}{RGB}{181,111,35}
\definecolor{CodeBg}{RGB}{248,248,248}

\lstdefinestyle{pythonstyle}{
  basicstyle=\ttfamily\small,
  backgroundcolor=\color{CodeBg},
  keywordstyle=\color{BDMBlue},
  commentstyle=\color{BDMGray},
  stringstyle=\color{BDMAmber},
  showstringspaces=false,
  breaklines=true,
  frame=single,
  rulecolor=\color{black!20},
  columns=fullflexible
}

\lstdefinestyle{sqlstyle}{
  language=SQL,
  basicstyle=\ttfamily\small,
  backgroundcolor=\color{CodeBg},
  keywordstyle=\color{BDMBlue},
  commentstyle=\color{BDMGray},
  stringstyle=\color{BDMAmber},
  showstringspaces=false,
  breaklines=true,
  frame=single,
  rulecolor=\color{black!20},
  columns=fullflexible
}

\lstset{style=pythonstyle}
\renewcommand{\lstlistingname}{代码}
\renewcommand{\bibname}{参考文献}

\theoremstyle{definition}
\newtheorem{definition}{定义}[chapter]
\newtheorem{example}[definition]{例}
\newtheorem{exercise}[definition]{习题}
\newtheorem{convention}[definition]{约定}

\theoremstyle{remark}
\newtheorem{remark}[definition]{评注}

\theoremstyle{plain}
\newtheorem{theorem}[definition]{定理}
\newtheorem{proposition}[definition]{命题}
\newtheorem{lemma}[definition]{引理}
\newtheorem{corollary}[definition]{推论}

\tcolorboxenvironment{definition}{
  colback=blue!2!white,
  colframe=blue!45!black,
  boxrule=0.4pt,
  arc=1.2mm,
  breakable
}

\tcolorboxenvironment{theorem}{
  colback=orange!3!white,
  colframe=orange!55!black,
  boxrule=0.4pt,
  arc=1.2mm,
  breakable
}

\tcolorboxenvironment{proposition}{
  colback=orange!2!white,
  colframe=orange!50!black,
  boxrule=0.4pt,
  arc=1.2mm,
  breakable
}

\tcolorboxenvironment{lemma}{
  colback=orange!2!white,
  colframe=orange!45!black,
  boxrule=0.4pt,
  arc=1.2mm,
  breakable
}

\tcolorboxenvironment{corollary}{
  colback=orange!2!white,
  colframe=orange!45!black,
  boxrule=0.4pt,
  arc=1.2mm,
  breakable
}

\newtcolorbox{chapterbox}{
  colback=gray!5!white,
  colframe=gray!55!black,
  boxrule=0.4pt,
  arc=1.2mm,
  breakable
}

\newtcolorbox{modernbox}[1][应用接口]{
  colback=green!3!white,
  colframe=green!45!black,
  boxrule=0.4pt,
  arc=1.2mm,
  title={#1},
  fonttitle=\mdseries,
  breakable
}

\newcommand{\R}{\mathbb R}
\newcommand{\N}{\mathbb N}
\newcommand{\E}{\mathbb E}
\newcommand{\Prob}{\mathbb P}
\newcommand{\dd}{\,\mathrm d}
\newcommand{\eps}{\varepsilon}
\newcommand{\abs}[1]{\left\lvert #1\right\rvert}
\newcommand{\norm}[1]{\left\lVert #1\right\rVert}
\newcommand{\argmin}{\operatorname*{arg\,min}}
\newcommand{\argmax}{\operatorname*{arg\,max}}


\title{《大数据与管理决策基础》\\[0.5em]\large 第2讲：企业数据与业务对象}
\author{课程讲义}
\date{}

\begin{document}

\maketitle
\tableofcontents

\setcounter{chapter}{1}
\chapter{企业数据与业务对象}
\label{ch:week02-enterprise-data-model}

\begin{chapterbox}
本讲讨论企业运营数据的建模基础。第一讲已经指出，数据分析的终点不是报表或模型本身，而是可执行、可评估、可反馈的管理行动。本讲进一步说明这些行动必须依赖稳定的数据对象：客户、订单、产品、库存、供应商、流程事件和业务状态。讲义从实体、属性、关系、事件、状态和指标等基本概念出发，介绍关系模型、键与完整性约束、ER 模型、规范化关系模型、星型模型、事实表粒度、度量可加性、时间建模和指标治理。通过本讲，学生应理解数据建模既是数据库设计问题，也是管理概念、组织责任和分析可信度的共同基础。
\end{chapterbox}

\begin{modernbox}[课程接口]
本讲把第一讲的管理行动闭环落到对象层。后续 SQL KPI、数据质量、统计看板、预测模型、优化决策、流程挖掘和 MLOps 都要回答同一个前置问题：模型和指标究竟作用于哪些业务对象，这些对象在数据表中如何被识别、连接、更新和审计。
\end{modernbox}

\section{学习目标}

第一讲讨论如何把经营症状转化为可行动的分析问题。第二讲回答一个更基础的问题：这些分析问题依赖什么样的数据结构？如果企业数据只是一些零散表格，那么后续 SQL、KPI、预测模型和优化模型都会缺少稳定基础。本讲的目标不是训练学生画一个漂亮的 ER 图，而是训练学生把管理概念转化为可查询、可治理、可建模、可行动的数据对象。

完成本讲后，学生应能够：
\begin{enumerate}
  \item 区分实体、属性、关系、事件、状态、事实和指标；
  \item 使用关系模式、主键、外键和完整性约束描述企业数据；
  \item 解释 ER 模型、规范化关系模型和星型模型各自适用的场景；
  \item 为订单、客户、产品和库存设计基本数据模型，并声明事实表粒度；
  \item 判断事实表度量的可加性，避免对比率、均值和快照余额进行错误聚合；
  \item 说明事件时间、状态时间和入库时间的差异，以及慢变维对历史分析的影响；
  \item 将关系表进一步映射为业务对象、对象关系、对象状态和可执行行动。
\end{enumerate}

\section{为什么管理决策需要数据建模}

在企业中，“有数据”并不等于“可决策”。同一个客户可能在 CRM、订单系统、售后系统和营销系统中拥有不同编号；同一笔销售可能被财务部门、运营部门和渠道部门以不同口径统计；同一件库存商品可能同时存在采购在途、仓库可用、销售锁定和质检冻结等状态。若没有明确的数据模型，后续分析会出现三类问题。

\begin{enumerate}
  \item 对象不稳定：系统无法判断同一个客户、订单或产品在不同系统中是否指向同一业务对象。
  \item 口径不稳定：不同报表对销售额、库存、延期、流失等指标使用不同定义。
  \item 行动不稳定：分析结果无法对应到可执行对象，例如无法说明对哪个订单加急、对哪个客户回访、对哪个库存位置补货。
\end{enumerate}

\begin{definition}[企业数据模型]
企业数据模型是对企业业务对象、对象属性、对象关系、业务事件、状态变化、指标口径和约束规则的结构化表达。它既可以表现为 ER 图、关系模式、维度模型，也可以表现为数据字典、指标字典和业务对象定义。
\end{definition}

这个定义强调，数据建模不是数据库工程师的局部工作，而是企业管理语言的工程化表达。没有数据模型，企业很难判断“客户”“订单”“库存”“销售额”“延期”这些词在不同系统和部门中是否具有同一含义。没有共同含义，就不可能有可信指标；没有可信指标，就不可能有可审计的决策。

\section{企业运营数据的基本语义}

企业运营数据看起来复杂，但从建模角度可以先分为六类。这个分类帮助学生判断一张表的角色：它是在描述一个对象，记录一次业务活动，保存某个时间点的状态，表达对象之间的连接，还是表达一个已经聚合过的指标。

\begin{longtable}{p{0.15\linewidth}p{0.32\linewidth}p{0.41\linewidth}}
\caption{企业运营数据的基本语义}\\
\toprule
类型 & 数据含义 & 示例\\
\midrule
\endfirsthead
\toprule
类型 & 数据含义 & 示例\\
\midrule
\endhead
实体 & 具有相对稳定身份的对象。 & 客户、产品、供应商、仓库、设备、员工。\\
属性 & 描述实体、关系或事件的变量。 & 地区、品类、单价、客户等级、付款方式。\\
关系 & 对象之间的结构性连接。 & 客户产生订单、产品属于品类、供应商供应物料。\\
事件 & 在时间上发生的一次业务活动。 & 下单、付款、发货、退货、维修、审批。\\
状态 & 某一时点或区间内对象的状态。 & 库存量、订单状态、设备状态、客户等级。\\
指标 & 对事件或状态的聚合度量。 & 销售额、准时交付率、库存周转率、流失率。\\
\bottomrule
\end{longtable}

\begin{definition}[业务事件]
业务事件是在特定时间发生、能够改变或记录业务状态的一次活动。事件通常具有事件主体、事件对象、发生时间、数量或金额等属性。
\end{definition}

一个订单明细事件可写为
\[
e_i=(o_i,c_i,p_i,q_i,t_i,s_i),
\]
其中 \(o_i\) 是订单，\(c_i\) 是客户，\(p_i\) 是产品，\(q_i\) 是数量，\(t_i\) 是事件时间，\(s_i\) 是交付或履约状态。这个表达式虽简单，却包含了企业数据建模的核心：客户和产品是相对稳定的实体，订单明细是交易事件，数量和时间是事件属性，交付状态可能来自后续物流事件或累计快照。

\begin{remark}
很多数据错误不是由计算公式造成的，而是由语义混淆造成的。例如，把库存余额当作销售事件求和，把客户当前等级用于解释历史订单，或把订单头表与订单明细表混在同一粒度上聚合，都会导致看似精确但实际错误的指标。
\end{remark}

\section{关系模型的基本语言}

Codd 提出的关系模型为数据库提供了严格而简洁的形式语言。在关系模型中，数据被表示为关系，关系可以理解为带有属性名的元组集合。一个关系模式可写为
\[
R(A_1:D_1,A_2:D_2,\ldots,A_k:D_k),
\]
其中 \(A_j\) 是属性名，\(D_j\) 是该属性的取值域。关系实例是满足该模式的一组元组：
\[
r(R)=\{t_1,t_2,\ldots,t_n\}.
\]
更严格地说，元组 \(t\) 可被看作从属性集合到相应取值域的映射：
\[
t(A_j)\in D_j,\quad j=1,\ldots,k.
\]

若客户集合为 \(C\)，产品集合为 \(P\)，时间集合为 \(T\)，订单明细集合可被看作有限关系：
\[
O \subseteq C\times P\times \mathbb N\times T\times S,
\]
其中 \(S\) 表示订单状态或履约状态的取值集合。

关系模型的重要优点是声明性。用户用 SQL 说明“需要什么结果”，而数据库系统决定如何执行连接、过滤、聚合和排序。这使企业数据可以被多个应用复用，也使管理分析能够建立在稳定的集合运算之上。但关系模型本身并不自动保证业务语义清楚。关系模式说明字段之间如何组织，业务模型还必须说明字段代表什么管理概念。

\section{键、外键与完整性约束}

在数据库设计中，至少需要关注三类约束。它们看起来像技术细节，但实际上直接影响管理分析的可信度。

\begin{enumerate}
  \item 实体完整性：主键不能为空且唯一。
  \item 引用完整性：外键必须引用已存在的主表记录。
  \item 业务完整性：数量、价格、日期等字段满足业务规则，例如数量为正、实际交付日期不早于下单日期。
\end{enumerate}

\begin{definition}[键]
给定关系 \(R\)，若属性集合 \(K\) 的取值能够唯一识别 \(R\) 中的每个元组，则 \(K\) 是超键。若 \(K\) 是超键且任意真子集都不是超键，则 \(K\) 是候选键。数据库设计者在候选键中选定用于标识元组的键称为主键。
\end{definition}

\begin{definition}[外键]
外键是一个关系中的属性集合，其取值引用另一个关系中的候选键或主键。外键把不同关系中的业务对象连接起来，是跨表分析和完整性检查的基础。
\end{definition}

\begin{table}[htbp]
\centering
\caption{自然键与代理键的比较}
\begin{tabular}{p{0.18\linewidth}p{0.34\linewidth}p{0.34\linewidth}}
\toprule
类型 & 优点 & 风险\\
\midrule
自然键 & 具有业务含义，便于人工理解。 & 业务规则变化时可能失效，例如手机号、证件号、商品编码变更。\\
代理键 & 稳定、短小，便于连接和仓库建模。 & 需要维护与业务标识之间的映射，不能替代业务唯一性校验。\\
\bottomrule
\end{tabular}
\end{table}

例如，如果订单表中的客户编号在客户表中不存在，那么按客户分层计算销售额时会出现无法归属的订单；如果同一个订单编号重复出现，则订单量、销售额和延期率都会被扭曲。因此，完整性约束不是数据库管理员的局部事务，而是指标可信度和模型可信度的前提。

\section{订单、客户、产品与库存的 ER 模型}

ER 模型强调从现实业务对象和对象关系出发建模，而不是直接从物理表结构出发建模。它尤其适合在课程早期训练学生识别实体、属性、关系、基数和可选性。

\begin{center}
\begin{tikzpicture}[
  entity/.style={draw, rounded corners=1pt, fill=blue!5, minimum width=2.8cm, minimum height=0.9cm, align=center},
  rel/.style={-{Stealth[length=2mm]}, thick}
]
\node[entity] (cust) {Customer\\客户};
\node[entity, right=3.1cm of cust] (order) {Order\\订单};
\node[entity, right=3.1cm of order] (prod) {Product\\产品};
\node[entity, below=1.6cm of order] (line) {OrderLine\\订单明细};
\node[entity, below=1.6cm of prod] (inv) {Inventory\\库存快照};
\draw[rel] (cust) -- node[above]{1:N} (order);
\draw[rel] (order) -- node[left]{1:N} (line);
\draw[rel] (prod) -- node[right]{1:N} (line);
\draw[rel] (prod) -- node[right]{1:N} (inv);
\end{tikzpicture}
\end{center}

该模型的关键点不是画图本身，而是明确若干会影响分析结果的建模选择：
\begin{itemize}
  \item 订单事实表的粒度是一行一张订单，还是一行一个订单明细项；
  \item 客户、产品等维度表是否能够稳定描述事实表中的对象；
  \item 库存是事件表、状态快照表，还是二者都需要；
  \item 每个关系是否能被主键和外键约束检查；
  \item 是否存在多对多关系，例如一个促销活动影响多个订单，一个订单使用多个优惠。
\end{itemize}

ER 模型一般经历概念模型、逻辑模型和物理模型三个层次。概念模型面向业务语言，强调对象和关系；逻辑模型转化为关系模式、键和约束；物理模型进一步考虑数据库类型、索引、分区、存储和性能。

\begin{table}[htbp]
\centering
\caption{三类模型层次}
\begin{tabular}{p{0.18\linewidth}p{0.31\linewidth}p{0.36\linewidth}}
\toprule
层次 & 主要问题 & 典型产物\\
\midrule
概念模型 & 业务上有哪些对象和关系？ & ER 图、对象定义、关系说明。\\
逻辑模型 & 如何转化为关系模式和约束？ & 表、字段、主键、外键、基数。\\
物理模型 & 如何在具体平台上实现？ & 索引、分区、存储格式、权限。\\
\bottomrule
\end{tabular}
\end{table}

\section{规范化与函数依赖}

规范化关系模型主要服务于一致性维护和事务处理。它通过减少冗余来避免插入异常、删除异常和更新异常。例如，若客户地址被重复写入每一笔订单，客户地址更新时就可能出现不同订单保存不同地址的矛盾。

规范化的基本思想可以用函数依赖表达。若属性集合 \(X\) 的取值确定属性集合 \(Y\) 的取值，则记为
\[
X\rightarrow Y.
\]
在客户表中，若客户编号确定客户名称和注册地区，则有
\[
\text{customer\_id}\rightarrow
\{\text{customer\_name},\text{registered\_region}\}.
\]

\begin{definition}[第三范式的直观要求]
一个关系模式接近第三范式，是指非键属性应直接依赖于键，而不应通过另一个非键属性间接依赖于键。其直观目标是减少冗余和更新异常。
\end{definition}

例如，在订单表中同时存储 \(\text{product\_id}\)、\(\text{product\_name}\) 和 \(\text{category}\) 通常会导致冗余，因为产品名称和品类应由产品编号决定，而不是由每一行订单重复决定。更合适的做法是将产品属性放入产品表，再通过 \(\text{product\_id}\) 与订单事实连接。

需要注意，规范化并不是越高越好。事务系统强调写入一致性，分析系统强调查询便利性和指标解释。二者的建模目标不同。

\section{事务型建模与分析型建模}

面向事务处理的系统通常需要高一致性、低延迟写入和明确业务流程；面向分析的系统通常需要跨主题查询、历史追踪和大量聚合。规范化模型与星型模型并不是谁更高级的关系，而是服务于不同目标。

\begin{table}[htbp]
\centering
\caption{事务型建模与分析型建模的差异}
\begin{tabular}{p{0.21\linewidth}p{0.31\linewidth}p{0.32\linewidth}}
\toprule
维度 & 事务型关系模型 & 分析型星型模型\\
\midrule
主要目标 & 支持写入、一致性和业务流程。 & 支持查询、聚合和指标分析。\\
冗余控制 & 尽量减少冗余。 & 允许适度冗余以简化查询。\\
典型结构 & 多张规范化业务表。 & 事实表加维度表。\\
时间处理 & 多关注当前业务状态。 & 强调历史快照和版本追踪。\\
关键风险 & 更新异常、引用错误。 & 粒度混乱、口径不一致。\\
\bottomrule
\end{tabular}
\end{table}

在管理分析中，常见做法是将交易事件建模为事实表，将描述对象的表建模为维度表。典型星型模型为
\[
\begin{aligned}
\text{FactOrderLine}\leftarrow
\{&\text{DimCustomer},\text{DimProduct},\text{DimDate},\\
&\text{DimChannel},\text{DimWarehouse}\}.
\end{aligned}
\]

\begin{definition}[事实表与维度表]
事实表记录可度量的业务事件或状态，通常包含外键、时间键和度量字段。维度表描述事实表中的分析角度，例如客户、产品、地区、渠道、日期和仓库。
\end{definition}

\section{事实表粒度}

事实表的设计首先要确定粒度（grain）。粒度指事实表中一行记录所代表的最小业务事实。例如：
\begin{itemize}
  \item 一行表示一个订单；
  \item 一行表示一个订单明细项；
  \item 一行表示一个产品、仓库、日期的库存快照；
  \item 一行表示一个客户、活动、日期的营销触达事件；
  \item 一行表示一个流程实例中的一个活动完成事件。
\end{itemize}

\begin{definition}[事实表粒度]
事实表粒度是事实表中每一行所代表的业务事实的最小单位。粒度必须在识别维度和度量之前声明，因为它决定哪些字段可以放入事实表，哪些指标可以被正确计算。
\end{definition}

Kimball 的四步维度建模方法可概括为：
\begin{enumerate}
  \item 选择业务过程；
  \item 声明事实表粒度；
  \item 识别维度；
  \item 识别事实或度量。
\end{enumerate}

第二步最容易被忽略。粒度决定了 KPI 的可计算性和可解释性。例如，销售额可定义为
\[
\text{Revenue}=\sum_i q_i\cdot price_i,
\]
但库存周转率需要同时使用销售成本和平均库存：
\[
\text{Inventory Turnover}
=
\frac{\text{Cost of Goods Sold}}{\text{Average Inventory}}.
\]
前者通常来自交易事实表，后者则需要库存状态快照或周期平均值。若把交易事件与状态快照混在一张事实表中，指标计算就会出现重复计数或时间口径错误。

\section{度量的可加性}

事实表中的度量并不都能用同一种方式聚合。可加性是指标建模中最容易被低估的概念。

\begin{table}[htbp]
\centering
\caption{事实表度量的可加性}
\begin{tabular}{p{0.18\linewidth}p{0.42\linewidth}p{0.28\linewidth}}
\toprule
类型 & 含义 & 示例\\
\midrule
可加度量 & 可沿所有维度求和。 & 销售金额、销售数量、成本。\\
半可加度量 & 只可沿部分维度求和。 & 库存余额可按产品求和，但不宜跨日期直接求和。\\
不可加度量 & 不能直接求和，通常需重新计算。 & 比率、百分比、均值、毛利率。\\
\bottomrule
\end{tabular}
\end{table}

\begin{proposition}[比率指标不能直接平均]
设第 \(d\) 日准时交付率为 \(r_d=m_d/n_d\)，其中 \(m_d\) 是准时交付订单数，\(n_d\) 是总订单数。一个周期的总体准时交付率应为
\[
\frac{\sum_d m_d}{\sum_d n_d},
\]
一般不等于每日准时交付率的简单平均
\[
\frac{1}{D}\sum_{d=1}^{D}r_d.
\]
\end{proposition}

该命题说明，指标字典不仅要写出指标名称，还要写出聚合规则。准时交付率、转化率、毛利率、流失率等比率指标都应尽量保留分子和分母，以便在不同维度和时间窗口上重新计算。

\section{事件、状态与时间}

企业系统中的时间数据至少有三种含义：事件发生时间、状态有效时间和数据入库时间。它们对应不同问题。

\begin{longtable}{p{0.20\linewidth}p{0.34\linewidth}p{0.34\linewidth}}
\caption{事件表与状态表的差异}\\
\toprule
表类型 & 适用问题 & 示例\\
\midrule
\endfirsthead
\toprule
表类型 & 适用问题 & 示例\\
\midrule
\endhead
交易事实表 & 记录一次业务活动。 & 一行订单明细、一笔付款、一条发货记录。\\
周期快照表 & 记录固定周期的状态。 & 每日产品、仓库库存余额。\\
累计快照表 & 跟踪有开始和结束的流程。 & 订单从下单、拣货、发货到签收的里程碑时间。\\
维度历史表 & 记录属性随时间变化。 & 客户等级、产品品类、组织归属的历史版本。\\
\bottomrule
\end{longtable}

慢变维（Slowly Changing Dimension, SCD）用于处理维度属性随时间变化的问题。例如客户从普通客户变为关键客户后，历史订单应按照下单当时的客户等级分析，还是按照当前客户等级分析？不同业务问题需要不同口径。因此，时间建模不是技术细节，而是企业指标解释的一部分。

\begin{center}
\captionsetup{type=table,hypcap=false}
\caption{慢变维的常见处理方式}
\begin{tabular}{p{0.22\linewidth}p{0.32\linewidth}p{0.34\linewidth}}
\toprule
方式 & 含义 & 适用问题\\
\midrule
覆盖更新 & 只保留最新属性值。 & 只关心当前状态的运营查询。\\
保留历史版本 & 通过生效时间和失效时间保存属性历史。 & 需要按历史口径解释订单或客户行为。\\
保留当前和历史字段 & 同时保留当前属性和部分历史属性。 & 需要兼顾当前管理视角和有限历史比较。\\
\bottomrule
\end{tabular}
\end{center}

\section{从关系表到业务对象}

关系表适合存储和查询，但管理行动通常作用于业务对象。一个订单对象可能来自订单表、订单明细表、客户表、产品表和物流状态表的组合；一个库存位置对象可能来自库存快照表、产品表、仓库表和锁定规则表的组合。

\begin{center}
\captionsetup{type=table,hypcap=false}
\caption{关系表到业务对象的映射}
\begin{tabular}{p{0.20\linewidth}p{0.40\linewidth}p{0.29\linewidth}}
\toprule
业务对象 & 来源表 & 可支持行动\\
\midrule
\texttt{Customer} & \texttt{customers}, \texttt{customer\_segments} & 分层运营、优惠券、客户回访。\\
\texttt{Product} & \texttt{products}, \texttt{product\_categories} & 定价、促销、补货策略。\\
\texttt{Order} & \texttt{orders}, \texttt{order\_lines}, \texttt{logistics} & 加急、拆单、客户沟通。\\
\texttt{InvPosition} & \texttt{inventory}, \texttt{products}, \texttt{warehouses} & 补货、调拨、清仓。\\
\texttt{Supplier} & \texttt{suppliers}, \texttt{purchase\_orders} & 供应商替换、交期协商。\\
\bottomrule
\end{tabular}
\end{center}

因此，现代企业决策系统关注
\[
\text{Tables}\rightarrow\text{Objects}\rightarrow\text{Links}\rightarrow\text{States}\rightarrow\text{Actions}.
\]

这一映射非常重要。SQL 查询返回的是行和列，而管理者执行的是对订单、客户、产品、库存位置等对象的行动。只有建立从表到对象的映射，后续的风险评分、权限控制、行动推荐和审计日志才有明确载体。

\section{指标字典、元数据与治理}

数据建模最终必须进入治理机制。DAMA-DMBOK 将数据治理、数据架构、数据建模与设计、数据质量、元数据、数据仓库、BI 和数据科学等看作相互关联的数据管理知识域。对于本课程而言，最重要的是理解三类文档。

\begin{enumerate}
  \item 数据字典：说明表、字段、类型、取值范围、主外键和业务含义。
  \item 指标字典：说明指标名称、业务定义、计算公式、数据来源、刷新频率、责任人和适用范围。
  \item 血缘说明：说明指标或模型特征从哪些源表经过哪些转换产生。
\end{enumerate}

\begin{table}[htbp]
\centering
\caption{指标字典的最小规范}
\begin{tabular}{p{0.23\linewidth}p{0.60\linewidth}}
\toprule
字段 & 说明\\
\midrule
指标名称 & 统一使用的业务名称。\\
业务定义 & 指标所刻画的经营含义。\\
计算公式 & 分子、分母、聚合规则、时间窗口和过滤条件。\\
数据来源 & 源表、字段、刷新频率和必要清洗规则。\\
粒度与维度 & 可以按哪些对象、时间和组织层级分析。\\
责任人 & 数据负责人、业务负责人和异常处理机制。\\
行动接口 & 指标变化后触发的预警、审批、复盘或业务动作。\\
\bottomrule
\end{tabular}
\end{table}

在企业项目中，数据字典和指标字典不是可有可无的附录。它们决定不同团队能否对同一事实达成一致，也决定模型结果是否可以追溯、审计和复现。

\section{课堂案例：零售订单数据模型}

考虑一家零售企业。第一讲中的经营症状是“库存金额持续上升，但热门商品仍然缺货”。若要进入数据分析，首先必须建立能够支持库存、销售和补货决策的数据模型。

\begin{table}[htbp]
\centering
\caption{零售库存问题的数据建模拆解}
\begin{tabular}{p{0.22\linewidth}p{0.62\linewidth}}
\toprule
建模要素 & 示例\\
\midrule
业务对象 & 客户、产品、仓库、订单、订单明细、库存位置。\\
交易事实 & 订单明细、付款、退货、补货入库、调拨出入库。\\
状态快照 & 每日产品、仓库库存余额，在途库存，锁定库存。\\
核心维度 & 日期、产品、品类、仓库、地区、渠道、客户分层。\\
关键指标 & 销售额、销量、缺货率、库存周转率、滞销库存金额。\\
可执行行动 & 补货、调拨、促销、清仓、供应商催交。\\
\bottomrule
\end{tabular}
\end{table}

在这个案例中，订单明细事实表回答“卖出了什么”，库存快照表回答“某天某仓库还剩多少”，补货事件表回答“何时补了多少”，产品维度表回答“这是什么品类和价格带”。如果把这些问题放在同一张宽表中，短期内可能方便查询，但长期会造成粒度混乱、历史状态丢失和指标解释困难。

\section{配套代码}

配套代码位于：
\begin{center}
\texttt{src/bdm\_decision/cases/week02\_enterprise\_data\_model.py}
\end{center}

代码的目的不是实现数据库系统，而是把本讲中的实体、字段、关系和业务对象映射转化为可执行的结构化表示。其核心抽象如下：
\begin{lstlisting}[language=Python,caption={实体、字段和关系}]
@dataclass(frozen=True)
class Entity:
    name: str
    primary_key: str
    fields: list[Field]
    description: str

@dataclass(frozen=True)
class Relationship:
    source: str
    target: str
    cardinality: str
    foreign_key: str
    description: str
\end{lstlisting}

运行方式：
\begin{lstlisting}[language=bash,caption={运行第二周案例代码}]
PYTHONPATH=src python3 src/bdm_decision/cases/week02_enterprise_data_model.py
\end{lstlisting}

学生阅读代码时应注意，\texttt{Entity} 和 \texttt{Relationship} 对应的是逻辑模型层，而 \texttt{business\_object\_mapping()} 对应的是从表到对象的管理接口。后续课程中的 SQL、指标和模型都可以接在这些对象之上。

\section*{习题}
\addcontentsline{toc}{section}{习题}

\begin{exercise}
请以“零售订单、客户、产品与库存”为背景完成：
\begin{enumerate}
  \item 列出至少四个实体及其主键；
  \item 为每个实体设计至少四个字段；
  \item 标明实体之间的基数关系；
  \item 指出一个交易事实表、一个周期快照表和至少三个维度表；
  \item 声明事实表粒度，并说明每个度量是否可加；
  \item 将至少两个关系表映射为业务对象，并说明其可执行行动；
  \item 为“销售额”“准时交付率”“库存周转率”写出指标口径。
\end{enumerate}
\end{exercise}

\begin{exercise}
某企业将每日转化率直接平均得到月度转化率。请构造一个两天的数值例子，说明这种做法可能产生误导，并给出正确计算方式。
\end{exercise}

\begin{exercise}
考虑客户等级从普通客户升级为关键客户的场景。请说明在分析历史销售额时，使用当前客户等级和使用历史客户等级会回答两个不同问题。
\end{exercise}

\section*{本讲小结}
\addcontentsline{toc}{section}{本讲小结}

本讲建立了企业数据建模的基础语言。实体、关系、事件、状态和指标是企业数据的基本语义单位；关系模型提供元组、属性、键、外键和完整性约束的形式基础；ER 模型帮助我们从业务对象出发讨论结构；规范化模型适合事务一致性，星型模型适合分析查询；事实表粒度、度量可加性和时间口径决定 KPI 是否可信；表到业务对象的映射则决定分析结果能否进入管理行动。后续课程的 SQL KPI、数据质量、看板、预测、优化和流程挖掘都将建立在这些概念之上。

\section*{常见误区}
\addcontentsline{toc}{section}{常见误区}

\begin{enumerate}
  \item 不区分订单头与订单明细，导致事实表粒度混乱；
  \item 使用业务含义不稳定的字段作为唯一主键；
  \item 只画 ER 图，不说明主外键和完整性约束；
  \item 将库存快照误认为交易事件；
  \item 在 KPI 计算前没有定义指标口径和事实表粒度；
  \item 对比率、百分比和均值进行简单相加或平均；
  \item 用当前维度属性解释历史事件，却没有说明时间口径；
  \item 把关系表直接等同于业务对象，忽略行动、权限和状态。
\end{enumerate}

\addcontentsline{toc}{chapter}{参考文献}
\begin{thebibliography}{99}
\bibitem{codd1970}
Codd, E. F. (1970).
\newblock A relational model of data for large shared data banks.
\newblock Communications of the ACM, 13(6), 377--387.

\bibitem{chen1976}
Chen, P. P. (1976).
\newblock The entity-relationship model: Toward a unified view of data.
\newblock ACM Transactions on Database Systems, 1(1), 9--36.

\bibitem{date2003}
Date, C. J. (2003).
\newblock An Introduction to Database Systems.
\newblock Addison-Wesley.

\bibitem{dama2017}
DAMA International. (2017).
\newblock DAMA-DMBOK: Data Management Body of Knowledge, 2nd edition.
\newblock Technics Publications.

\bibitem{golfarelli2009}
Golfarelli, M., and Rizzi, S. (2009).
\newblock Data Warehouse Design: Modern Principles and Methodologies.
\newblock McGraw-Hill.

\bibitem{inmon2005}
Inmon, W. H. (2005).
\newblock Building the Data Warehouse.
\newblock Wiley.

\bibitem{kimball2013}
Kimball, R., and Ross, M. (2013).
\newblock The Data Warehouse Toolkit: The Definitive Guide to Dimensional Modeling.
\newblock Wiley.
\end{thebibliography}

\end{document}
